\origpage{357}{N. Robertson, P.D. Seymour / Journal of Combinatorial Theory, Series B 92 (2004) 325--357}
\textit{Suppose that for each $i \ge 1$ and each tangle $\mathcal{T}$ in $G_i$ of order $\ge \theta$ which controls a $K_p$-minor of $sk(G_i^-)$, there is a rooted location $\mathcal{L}$ in $G_i$ which $\theta$-isolates $\mathcal{T}$ such that $(P_i, \mathcal{L})$ has a heart in $\mathcal{F}$. Then there exist $j > i \ge 1$ such that $P_i$ is simulated in $P_j$.}

\textbf{Proof.} Choose $\theta_1$ and $\mathcal{F}_1$ such that 11.1 holds (with $\theta$, $\mathcal{F}$ replaced by $\theta_1$, $\mathcal{F}_1$). Let $\mathcal{F}_2$ be the set of partial $\Omega$-patchworks $(G, \mu, \Delta, \phi)$ with $dom(\phi) = \emptyset$ and $(G, \mu, \Delta) \in \mathcal{F}_1$. Then $\mathcal{F}_2$ is well-behaved. Let $\mathcal{F}_3 = \mathcal{F} \cup \mathcal{F}_2$; then $\mathcal{F}_3$ is well-behaved. Moreover, for each $i \ge 1$ and each tangle $\mathcal{T}$ in $G_i$ of order $\ge \theta_2 = \max(\theta, \theta_1)$, there is a rooted location $\mathcal{L}$ in $G_i$ such that $\mathcal{L}$ $\theta_2$-isolates $\mathcal{T}$ and $(P_i, \mathcal{L})$ has a heart in $\mathcal{F}_3$; for if $\mathcal{T}$ controls a $K_p$-minor of $sk(G_i^-)$, this is true by hypothesis, and if not then this is true by 11.1. The result follows from 3.1. $\square$

\sectionhead{Acknowledgment}

We thank the referees for their diligence with this paper, and for their very helpful comments.

\sectionhead{References}

[1] N. Robertson, P.D. Seymour, Graph minors. IV. Tree-width and well-quasi-ordering, J. Combin. Theory Ser. B 48 (1990) 227--254. [2] N. Robertson, P.D. Seymour, Graph minors. X. Obstructions to tree-decomposition, J. Combin. Theory Ser. B 52 (1991) 153--190. [3] N. Robertson, P.D. Seymour, Graph minors. XVII. Taming a vortex, J. Combin. Theory Ser. B 77 (1999) 162--210. [4] N. Robertson, P.D. Seymour, Graph minors. XVIII. Tree-decompositions and well-quasi-ordering, J. Combin. Theory Ser. B 89 (2003) 77--108. [5] N. Robertson, P.D. Seymour, Graph minors. XIX. Well-quasi-ordering on a surface, J. Combin. Theory Ser. B 90 (2004) 325--385. [6] K. Wagner, Graphentheorie, vol. 248/248a, B. J. Hochschultaschenbucher, Mannheim, 1970, p. 61.
